\begin{figure*}
\hspace{-1cm}
\begin{subfigure}[t]{0.43\linewidth}
\vspace{0pt}
%% This code snippet is in a table for 2 reasons: to center it on the page,
%% and to draw a tight bounding box around it.
\begin{tabular}{|c|c} 
\hline
\hspace{2mm}
{
\lstset{numbers=left,frame=l,numbersep=6pt,numberstyle=\color{black}}
\begin{code}
data Exp = Lit Int
         | Plus Exp Exp
         | Sub Exp Exp
         | Let Sym Exp Exp
         ...

constFold :: Exp -> Exp
constFold exp = case exp of
  Lit i -> Lit i
  Plus e1 e2 ->
    case (e1, e2) of
      (Lit i,Lit j) -> Lit (i+j)
      _ -> let (e3,e4) =
             ( constFold e1 @\partup@ @\label{line:parallel-tuple}@
               constFold e2 )
           in Plus e3 e4
  Sub e1 e2 ->
    ...
\end{code}
}
%% \hspace{3mm}
\\
\hline
\end{tabular}
\caption{Constant folding written using the front-end language for Gibbon (Haskell).}
\label{fig:example1-haskell}
\end{subfigure}
%
\hspace{0.04\linewidth}
%
\begin{subfigure}[t]{0.43\linewidth}
\vspace{0pt}
%% %% This code snippet is in a table for 2 reasons: to center it on the page,
%% %% and to draw a tight bounding box around it.
%% \begin{figure}[h!]
\begin{tabular}{|c|c}
\hline
\hspace{2mm}
{
\lstset{numbers=left,frame=l,numbersep=6pt,numberstyle=\color{black}}
\begin{code}
constFold: forall@\locreg{l_1}{r_1}@ @\locreg{l_2}{r_2}@. @\styatlocreg{Exp}{l_1}{r_1}@ -> @\styatlocreg{Exp}{l_2}{r_2}@
constFold [@\locreg{l_1}{r_1}@ @\locreg{l_2}{r_2}@] exp = case exp of
  Lit (i:@\styatlocreg{Int}{l_i}{r_1}@) -> (Lit @\locreg{l_2}{r_2}@ i) @\label{line:datacon}@
  Plus (e1: @\styatlocreg{Exp}{l_a}{r_1}@) (e2: @\styatlocreg{Exp}{l_b}{r_1}@) ->
    case (e1,e2) of
      (Lit(i:@\styatlocreg{Int}{l_c}{r_1}@), Lit(j:@\styatlocreg{Int}{l_d}{r_1}@))
        -> (Lit @\locreg{l_2}{r_2}@ (i+j))
      _
        ->
        letloc @\locreg{l_3}{r_2}@ = @\locreg{l_2}{r_2}@ + 1 in @\label{line:letloctag}@
        let e3 : @\styatlocreg{Exp}{l_3}{r_2}@ =
          constFold [@\locreg{l_a}{r_1}@ @\locreg{l_3}{r_2}@] e1 in
        letloc @\locreg{l_4}{r_2}@ = after(@\styatlocreg{Exp}{l_3}{r_2}@) in @\label{line:letlocafter}@
        let e4 : @\styatlocreg{Exp}{l_4}{r_2}@ =
          constFold [@\locreg{l_b}{r_1}@ @\locreg{l_4}{r_2}@] e2 in
        (Plus @\locreg{l_2}{r_2}@ e3 e4) @\label{line:joinpoint}@
  Sub (e1: @\styatlocreg{Exp}{l_a}{r_1}@) (e2: @\styatlocreg{Exp}{l_b}{r_1}@) ->
    ...
\end{code}
}
\hspace{1cm}
\\
\hline
\end{tabular}
\caption{\Figref{fig:example1-haskell} compiled into LoCal IR by Gibbon.}
\label{fig:example1}
\end{subfigure}
\vspace{-5mm}
\caption{\vspace{-3mm}}
%% \caption{}
%% \caption{{\bf Left (a)}: Constant folding written using the front-end language for Gibbon (Haskell).
%% {\bf Right (b)}: the same program compiled into LoCal IR by Gibbon.}
\label{fig:example1-both}
\end{figure*}
